\section{Relation to existing M\"obius correlation theorems}
\label{sec:tpc41-literature}

The scalar theorem in \cref{sec:tpc41-triple-prime} uses a result designed
for large moduli.  Other major advances on M\"obius correlations illuminate
the boundary but do not directly estimate the physical form.

\subsection{Short arithmetic progressions}

The theorem of \citet{KlurmanMangerelTeravainen2023} gives, outside a
nearly power--saving exceptional set of moduli, centered variance
\begin{equation}
 \ll\eps\varphi(M)(X/M)^2.
 \label{eq:tpc41-literature-KMT}
\end{equation}
At the TPC scale this is \(\eps XK\).  Our
\cref{lem:tpc41-character-disposal} removes their rank--one term for
almost all members of the declared triple--prime family, so that the raw
M\"obius vector has the same logarithmic saving.

The scalar saving by itself is not new.  The classical
Barban--Davenport--Halberstam theorem for \(\mu\), as recalled in
\citet[Section~3.1.3]{KlurmanMangerelTeravainen2023}, gives for every
fixed \(A>0\), with a suitable \(B=B(A)\),
\begin{equation}
 \sum_{q\le X/(\log X)^B}
 \sum_{a\bmod q}^{*}
 \left|\sum_{\substack{n\le X\\n\equiv a\pmod q}}\mu(n)\right|^2
 \ll_A\frac{X^2}{(\log X)^A}.
 \label{eq:tpc41-literature-BDH}
\end{equation}
Because \(M_0=X^{399/400}\) eventually lies below every such cutoff,
restriction to \(\mathfrak M_X\) and Markov's inequality already give
arbitrary fixed logarithmic savings outside a logarithmic--density
exception.  The additional content of our KMT--large--sieve route is the
explicit stronger exceptional proportion
\eqref{eq:tpc41-raw-exception-proportion}, together with the admissible
three--interval and bounded--variation terminal specialization.

The result remains above the scalar atomic diagonal:
\begin{equation}
 \frac{\eps XK}{X}=\eps K
 =\frac{X^{1/400}}{(\log X)^\theta}\longrightarrow\infty.
 \label{eq:tpc41-literature-missing-K}
\end{equation}
Thus it is not a diagonal--scale progression theorem.  Assuming GRH can
remove or reduce exceptional--modulus issues in the cited framework, but
it does not by itself supply the missing factor
\(K/(\log X)^\theta\) in
\eqref{eq:tpc41-literature-missing-K}.

\subsection{Averaged Chowla and fixed affine forms}

The averaged Chowla theorem of
\citet{MatomakiRadziwillTao2015} gains cancellation after averaging over
all shifts in a dense box.  The alias differences here are the sparse
multiples \(rM\), and the physical cross--row kernel also has growing
slopes \(m_\alpha,m_\beta\).  Restricting a dense shift average to that
sparse subset loses the density factor that one is trying to save.

Results at almost all scales for fixed shifts
\citep{TaoTeravainen2019} and logarithmically averaged two--point results
\citep{Tao2016} likewise do not give a uniform estimate for
\eqref{eq:tpc41-affine-forms} when all coefficients vary with \(X\).
The Bohr--set analogue of \citet{TeravainenWalker2025} has fixed Bohr
data and logarithmic averaging; the TPC modulus, frequencies, row slopes,
and masks all move.

\subsection{Shifted primes}

The ideal convolution shadow in
\eqref{eq:tpc41-completed-shifted-prime} lies within the averaged-shift
theory of \citet{Lichtman2022,LichtmanTeravainen2022}.  This is a genuine
positive indication for the completion route.  It does not cover a fixed
twin shift with the actual localized vector weights.  Averaging over
\(h\) would change the problem: it may prove an almost--all--shift theorem,
but it cannot by itself close the prescribed fixed--\(h\) physical packet.

\subsection{The parity boundary}

The surviving form contains an even number of M\"obius factors and is
attached to a prime--pair sieve endpoint.  This is consistent with the
classical parity obstruction \citep{HeathBrown1982Parity}.  The present
paper makes two limited advances:
\begin{enumerate}[label=(\roman*)]
 \item it proves that the same--row two--M\"obius alias sector was a Gram
 artifact and removes it unconditionally; and
 \item it obtains a logarithmic saving for the raw scalar CRT residue
 vector on almost all admissible triple--prime moduli.
\end{enumerate}
Neither statement estimates the cross--row form
\eqref{eq:tpc41-physical-FL}.  We therefore do not describe either one as
a breach of parity.
